\uuid{ekRr}
\titre{Quotient/Composée/DL usuels}
\auteur{hanani ; mzari}
\organisation{exo7}
\licence{CC BY-NC-SA 4.0}
\datecreate{2020-01-20}
\niveau{L1}
\module{Analyse}
\chapitre{Développement limité}
\sousChapitre{Calculs}
\difficulte{3}
\theme{Opérations sur les DL, Quotient, Composée, DL usuels}
\exo7qcmid{lille-2.developpements-limites.026}

\contenu{
\texte{
Soit $a\in \Rr$ et $f$ une fonction. On dit que $f(x)$ est un $o\left((x-a)^n\right)$ au voisinage de $a$ et on écrit $f(x)=o((x-a)^n)$ si $\displaystyle \lim _{x\to a}\frac{f(x)}{(x-a)^n}=0$. On note par $DL_n(a)f(x)$ le développement limité de $f$ à l'ordre $n$ en $a$. On note aussi $DL_n(a^+)f(x)$ (resp. $DL_n(a^-)f(x)$) le développement limité de $f$ à l'ordre $n$ à droite en $a$ (resp. à gauche en $a$).
}

\begin{question}
\qclassification{125.02}
\qtags{motcle=Quotient/Composée/DL usuels}
Soit $\displaystyle f(x)=(1+x)^{1/(\mathrm{e}^x-1)}$. Parmi les affirmations suivantes, cocher celles qui sont vraies :
\begin{answers}
\good{$\displaystyle DL_2(0)\ln (1+x)=x-\frac{x^2}{2}+o(x^2)$ et $\displaystyle DL_2(0)(\mathrm{e}^x-1)=x+\frac{x^2}{2}+o(x^2)$.}
\bad{$\displaystyle DL_2(0)\frac{\ln (1+x)}{\mathrm{e}^x-1}=1-x+\frac{x^2}{2}+o(x^2)$.}
\bad{$\displaystyle DL_2(0)f(x)=\mathrm{e}.\mathrm{e}^{-x+x^2/2+o(x^2)}=\mathrm{e}-\mathrm{e}.x+\mathrm{e}.x^2+o(x^2)$.}
\good{$\displaystyle DL_2(0)f(x)=\mathrm{e}-\mathrm{e}.x+\frac{7\mathrm{e}}{6}x^2+o(x^2)$.}
\end{answers}
\vskip2mm
\begin{explanations}
On a : $\displaystyle f(x)=\mathrm{exp}\left(\frac{\ln (1+x)}{\mathrm{e}^x-1}\right)$. Donc, pour écrire le $\displaystyle DL_2(0)f(x)$, on utilise le $DL_3(0)\ln (1+x)$ et le $DL_3(0)(\mathrm{e}^x-1)$. La division suivants les puissances croissantes donne
$$DL_2(0)\frac{\ln (1+x)}{\mathrm{e}^x-1}=1-x+\frac{2x^2}{3}+o(x^2).$$
Ensuite, avec $\displaystyle \mathrm{e}^{1+u}=\mathrm{e}.\mathrm{e}^u=\mathrm{e}.\left(1+u+\frac{u^2}{2}+o(u^2)\right)$ où $\displaystyle u=-x+\frac{2x^2}{3}+o(x^2)$, on obtient :
$$f(x)=\mathrm{e}-\mathrm{e}.x+\frac{7\mathrm{e}}{6}x^2+o(x^2).$$
\end{explanations}
\end{question}
}
